\chapter{آشنایی با رباتی که می‌سازیم}

\begin{goalbox}
نتیجهٔ نهایی، ابزارهای اندک مورد نیاز و ساختار پروژه را بشناسید.
\end{goalbox}

هدف ما آگاهانه کوچک انتخاب شده است: ساخت ربات تلگرامی که تمام‌کردن و فهمیدن
آن آسان باشد. این ربات کارهای زیر را انجام می‌دهد:

\begin{itemize}
\item وقتی کاربر دستور
\lr{/start}
را می‌فرستد، به او خوشامد می‌گوید؛
\item با دریافت دستور
\lr{/help}
روش کار خود را توضیح می‌دهد؛
\item به متن عادی با عبارت
\lr{You said:}
و سپس همان متن پاسخ می‌دهد؛
\item پس از دریافت عکس، پیام صوتی یا محتوای مشابه، از کاربر متن می‌خواهد.
\end{itemize}

برای نمونه، وقتی علیرضا پیام زیر را می‌فرستد، گفت‌وگو چنین خواهد بود:

\begin{LTR}
\begin{resultbox}
\textbf{Alireza:} My first bot\\
\textbf{Bot:} You said: My first bot
\end{resultbox}
\end{LTR}

\section{چهار بخش اصلی}

تلگرام گفت‌وگو را نمایش می‌دهد.
\lr{BotFather}
حساب ربات را می‌سازد و توکن محرمانهٔ آن را تحویل می‌دهد. پایتون دستورهای ما را
اجرا می‌کند. بستهٔ
\lr{pyTelegramBotAPI}
ماژول
\lr{telebot}
را فراهم می‌کند تا این دستورها به تلگرام متصل شوند. در این آموزش از روش
\lr{polling}
استفاده می‌کنیم؛ یعنی تا وقتی برنامه باز است، مرتب از تلگرام می‌پرسد آیا پیام
تازه‌ای رسیده است یا نه.

\section{ساختاری کوچک و خوانا}

\begin{LTR}
\begin{lstlisting}[language={},numbers=none]
telbot-simple-bot/
|-- code/
|   |-- bot.py
|   |-- requirements.txt
|   |-- test_bot.py
|   `-- README.md
|-- documentation/
|   |-- English/
|   `-- Persian/
|-- LINKEDIN_POST.md
`-- README.md
\end{lstlisting}
\end{LTR}

پس از نصب وابستگی، فقط فایل
\lr{code/bot.py}
برای اجرای ربات لازم است. فایل‌های دیگر برای آموزش، آزمون یا مستندسازی هستند.

\begin{checkpointbox}
اکنون دقیقاً می‌دانید ربات نهایی چه می‌کند. هنوز هیچ توکن، بسته یا حسابی تغییر
نکرده است.
\end{checkpointbox}
